\section{Worked Examples}

\bigskip

\subsection{PCA with Two Variables}
\label{worked-pca-two-variables}

This example illustrates the transition from original variables to principal components using a bivariate dataset.

\medskip

\begin{enumerate}
    \setlength{\itemsep}{1em}
    
    \item \textit{Original Data}: Suppose we have the following observations for variables $X_1$ and $X_2$:

    \[
    \mathbf{X} =
    \begin{bmatrix}
    2 & 3 \\
    4 & 5 \\
    5 & 7 \\
    3 & 4
    \end{bmatrix}
    \]

    \item \textit{Standardizing the Data}: Standardize each variable to have mean 0 and standard deviation 1:

    \[
    Z_1 = \frac{X_1 - \bar{X}_1}{s_{X_1}}, \quad
    Z_2 = \frac{X_2 - \bar{X}_2}{s_{X_2}}
    \]

    Compute means and standard deviations:

    \[
    \bar{X}_1 = \frac{2+4+5+3}{4} = 3.5, \quad
    \bar{X}_2 = \frac{3+5+7+4}{4} = 4.75
    \]
    
    \[
    s_{X_1} = \sqrt{\frac{(2-3.5)^2 + (4-3.5)^2 + (5-3.5)^2 + (3-3.5)^2}{4-1}} \approx 1.29
    \]
    
    \[
    s_{X_2} = \sqrt{\frac{(3-4.75)^2 + (5-4.75)^2 + (7-4.75)^2 + (4-4.75)^2}{3}} \approx 1.71
    \]

    Standardized matrix:

    \[
    \mathbf{Z} =
    \begin{bmatrix}
    -1.16 & -1.03 \\
     0.39 & 0.15 \\
     1.16 & 1.31 \\
    -0.39 & -0.22
    \end{bmatrix}
    \]

    Total variance of standardized variables: $1+1=2$.

    \item \textit{Correlation Matrix}: Compute the correlation (equivalent to covariance for standardized data):
    
    The correlation between $Z_1$ and $Z_2$ is
    
    \[
    r_{12} = \frac{1}{n-1} \sum_{i=1}^{n} Z_{1i} Z_{2i}
    \]
    
    Step by step multiplication:
    
    \[
    Z_1 Z_2 =
    \begin{bmatrix}
    (-1.16)(-1.03) = 1.1948 \\
    (0.39)(0.15) = 0.0585 \\
    (1.16)(1.31) = 1.5196 \\
    (-0.39)(-0.43) = 0.1677
    \end{bmatrix}
    \]
    
    Sum:
    
    \[
    1.1948 + 0.0585 + 1.5196 + 0.1677 \approx 2.9406
    \]
    
    Divide by $n-1 = 3$:
    
    \[
    r_{12} = \frac{2.9406}{3} \approx 0.9802
    \]
    
    Correlations of variables with themselves are 1, so the correlation matrix is
    
    \[
    \mathbf{R} =
    \begin{bmatrix}
    1 & 0.980 \\
    0.980 & 1
    \end{bmatrix}
    \]
    
    \item \textit{Eigenvalues and Eigenvectors}:
    
    Solve the characteristic equation
    
    \[
    \det(\mathbf{R} - \lambda \mathbf{I}) = 0
    \]
    
    Step by step:
    
    \[
    \begin{vmatrix}
    1-\lambda & 0.980 \\
    0.980 & 1-\lambda
    \end{vmatrix} = (1-\lambda)(1-\lambda) - (0.980)^2 = 0
    \]
    
    \[
    (1-\lambda)^2 - 0.9604 = 0
    \]
    
    \[
    (1-\lambda)^2 = 0.9604 \implies 1-\lambda = \pm 0.980
    \]
    
    \[
    \lambda_1 = 1 + 0.980 = 1.980, \quad \lambda_2 = 1 - 0.980 = 0.020
    \]
    
    Eigenvectors:
    
    For $\lambda_1 = 1.980$:
    
    \[
    (\mathbf{R} - \lambda_1 \mathbf{I})\mathbf{v}_1 =
    \begin{bmatrix}
    1-1.980 & 0.980 \\
    0.980 & 1-1.980
    \end{bmatrix} 
    \begin{bmatrix} v_{11} \\ v_{12} \end{bmatrix} =
    \begin{bmatrix}
    -0.980 & 0.980 \\
    0.980 & -0.980
    \end{bmatrix}
    \begin{bmatrix} v_{11} \\ v_{12} \end{bmatrix} = 0
    \]
    
    Solving $-0.980 v_{11} + 0.980 v_{12} = 0 \implies v_{11} = v_{12}$. Normalize:
    
    \[
    \mathbf{v}_1 = \frac{1}{\sqrt{2}}
    \begin{bmatrix} 1 \\ 1 \end{bmatrix} =
    \begin{bmatrix} 0.707 \\ 0.707 \end{bmatrix}
    \]
    
    For $\lambda_2 = 0.020$:
    
    \[
    v_{21} = -v_{22} \implies
    \mathbf{v}_2 = \frac{1}{\sqrt{2}}
    \begin{bmatrix} 1 \\ -1 \end{bmatrix} =
    \begin{bmatrix} 0.707 \\ -0.707 \end{bmatrix}
    \]


    Eigenvectors (normalized):

    \[
    \mathbf{v}_1 = 
    \begin{bmatrix}
    0.707 \\
    0.707
    \end{bmatrix}, \quad
    \mathbf{v}_2 =
    \begin{bmatrix}
    0.707 \\
    -0.707
    \end{bmatrix}
    \]

    \item \textit{Forming Principal Components}: Each component is a linear combination of standardized variables:

    \[
    \xi_1 = 0.707 Z_1 + 0.707 Z_2, \quad
    \xi_2 = 0.707 Z_1 - 0.707 Z_2
    \]

    The first component explains $\frac{1.981}{2} \approx 99\%$ of the total variance.

    \item \textit{Verifying Orthonormality}:

    \[
    0.707^2 + 0.707^2 \approx 0.5 + 0.5 = 1
    \]

    \item \textit{Calculating Component Scores} (coordinates of the original data in the new PC space):

    \[
    \mathbf{\Xi} = \mathbf{Z} \mathbf{V} =
    \begin{bmatrix}
    -1.46 & -0.09 \\
     0.38 & -0.04 \\
     1.82 & 0.07 \\
    -0.74 & 0.05
    \end{bmatrix}
    \]

    Here, the first column is $\xi_1$, the second is $\xi_2$.

    \item \textit{Reconstructing Original Variables}:

    Using the first component only:

    \[
    Z_1 \approx 0.707 \xi_1, \quad Z_2 \approx 0.707 \xi_1
    \]

    Including both components:

    \[
    \mathbf{Z} = \xi_1 \mathbf{v}_1^\top + \xi_2 \mathbf{v}_2^\top
    \]

\end{enumerate}



\newpage


\subsection{MDS Example}
\label{worked-mds}

\medskip

Consider four objects with distance matrix
\[
\mathbf{\Delta}
=
\begin{pmatrix}
0 & 3 & 4 & 5 \\
3 & 0 & 5 & 6 \\
4 & 5 & 0 & 7 \\
5 & 6 & 7 & 0
\end{pmatrix}.
\]

\medskip
\noindent
\textbf{Step 1: Square the distances}

\medskip

Squaring each entry gives
\[
\mathbf{\Delta}^{(2)}
=
\begin{pmatrix}
0 & 9 & 16 & 25 \\
9 & 0 & 25 & 36 \\
16 & 25 & 0 & 49 \\
25 & 36 & 49 & 0
\end{pmatrix}.
\]

\medskip
\noindent
\textbf{Step 2: Construct the centering matrix}

\medskip

Since $n = 4$,
\[
\mathbf{J}
=
\mathbf{I}_4 - \frac{1}{4}\mathbf{1}\mathbf{1}^T.
\]

Because $\frac{1}{4}\mathbf{1}\mathbf{1}^T$ is a matrix of all $1/4$'s, we obtain
{
\renewcommand{\arraystretch}{1.5}
\[
\mathbf{J}
=
\begin{pmatrix}
\frac{3}{4} & -\frac{1}{4} & -\frac{1}{4} & -\frac{1}{4} \\
-\frac{1}{4} & \frac{3}{4} & -\frac{1}{4} & -\frac{1}{4} \\
-\frac{1}{4} & -\frac{1}{4} & \frac{3}{4} & -\frac{1}{4} \\
-\frac{1}{4} & -\frac{1}{4} & -\frac{1}{4} & \frac{3}{4}
\end{pmatrix}.
\]}

\medskip
\noindent
\textbf{Step 3: Compute row means, column means, and the grand mean}

\medskip

Row means of $\mathbf{\Delta}^{(2)}$:

\[
\bar{a}_{1\cdot} = \frac{0+9+16+25}{4} = 12.5
\]
\[
\bar{a}_{2\cdot} = \frac{9+0+25+36}{4} = 17.5
\]
\[
\bar{a}_{3\cdot} = \frac{16+25+0+49}{4} = 22.5
\]
\[
\bar{a}_{4\cdot} = \frac{25+36+49+0}{4} = 27.5
\]

Because the matrix is symmetric, the column means are the same.

The grand mean is
\[
\bar{a}_{\cdot\cdot}
=
\frac{1}{16}
\sum_{i,j} a_{ij}
=
\frac{80}{4}
=
20.
\]

\medskip
\noindent
\textbf{Step 4: Apply double centering and rescaling}

\medskip

Recall that
\[
b_{ij}
=
-\frac{1}{2}
\left(
a_{ij}
-
\bar{a}_{i\cdot}
-
\bar{a}_{\cdot j}
+
\bar{a}_{\cdot\cdot}
\right).
\]

We compute a few entries explicitly.

\medskip

Entry $(1,1)$:
\[
b_{11}
=
-\frac{1}{2}
\left(
0 - 12.5 - 12.5 + 20
\right)
=
-\frac{1}{2}(-5)
=
2.5.
\]

\medskip

Entry $(1,2)$:
\[
b_{12}
=
-\frac{1}{2}
\left(
9 - 12.5 - 17.5 + 20
\right)
=
-\frac{1}{2}(-1)
=
0.5.
\]

\medskip

Entry $(1,4)$:
\[
b_{14}
=
-\frac{1}{2}
\left(
25 - 12.5 - 27.5 + 20
\right)
=
-\frac{1}{2}(5)
=
-2.5.
\]

Carrying this out for all entries yields
\[
\mathbf{B}
=
\begin{pmatrix}
2.5 & 0.5 & -0.5 & -2.5 \\
0.5 & 7.5 & -2.5 & -5.5 \\
-0.5 & -2.5 & 12.5 & -9.5 \\
-2.5 & -5.5 & -9.5 & 17.5
\end{pmatrix}.
\]

\medskip
\noindent
\textbf{Step 5: Full eigen-decomposition of $\mathbf{B}$}

\medskip

The eigenvalues of $\mathbf{B}$ are approximately
\[
\lambda_1 = 25.4573, \quad
\lambda_2 = 11.6406, \quad
\lambda_3 = 2.9021, \quad
\lambda_4 \approx 0.
\]

Thus the diagonal matrix of eigenvalues is

\[
\mathbf{\Lambda}
=
\begin{pmatrix}
25.4573 & 0 & 0 & 0 \\
0 & 11.6406 & 0 & 0 \\
0 & 0 & 2.9021 & 0 \\
0 & 0 & 0 & 0
\end{pmatrix}.
\]

A corresponding orthonormal eigenvector matrix is

\[
\mathbf{V}
=
\begin{pmatrix}
-0.0621 & 0.1435 & -0.9445 & 0.2887 \\
-0.2334 & 0.8202 & 0.3365 & -0.4082 \\
-0.5738 & -0.5204 & 0.0031 & -0.5774 \\
0.7833 & -0.1869 & 0.0008 & -0.5774
\end{pmatrix}.
\]

\medskip
\noindent
\textbf{Step 6: Constructing the 2D configuration}

\medskip

We retain the first two eigenvalues and eigenvectors:

\[
\mathbf{\Lambda}_2
=
\begin{pmatrix}
25.4573 & 0 \\
0 & 11.6406
\end{pmatrix},
\qquad
\mathbf{V}_2
=
\begin{pmatrix}
-0.0621 & 0.1435 \\
-0.2334 & 0.8202 \\
-0.5738 & -0.5204 \\
0.7833 & -0.1869
\end{pmatrix}.
\]

Taking square roots of the eigenvalues,

\[
\mathbf{\Lambda}_2^{1/2}
=
\begin{pmatrix}
\sqrt{25.4573} & 0 \\
0 & \sqrt{11.6406}
\end{pmatrix}
=
\begin{pmatrix}
5.0455 & 0 \\
0 & 3.4118
\end{pmatrix}.
\]

The coordinate matrix is

\[
\mathbf{X}
=
\mathbf{V}_2 \mathbf{\Lambda}_2^{1/2}.
\]

Carrying out the multiplication yields

\[
\mathbf{X}
=
\begin{pmatrix}
-0.402 & 0.490 \\
-0.870 & 2.490 \\
-2.808 & -2.108 \\
4.081 & -0.872
\end{pmatrix}.
\]

\medskip
\noindent
\textbf{Recomputing Euclidean distances}

\medskip

We now verify that these coordinates reproduce the original distances.

\medskip

Distance between objects 1 and 2:
\[
\sqrt{
(-0.402 + 0.870)^2
+
(0.490 - 2.490)^2
}
=
\sqrt{
(0.468)^2 + (-2.000)^2
}
=
\sqrt{0.219 + 4}
=
\sqrt{4.219}
\approx 2.054.
\]

\medskip

Distance between objects 1 and 3:
\[
\sqrt{
(-0.402 + 2.808)^2
+
(0.490 + 2.108)^2
}
=
\sqrt{
(2.406)^2 + (2.598)^2
}
=
\sqrt{5.791 + 6.752}
=
\sqrt{12.543}
\approx 3.543.
\]
This does not exactly match the original distance of 4.
The discrepancy is not merely rounding error. It arises because we
retained only the first two eigenvalues when constructing the
configuration.

Recall that $\mathbf{B}$ had a third positive eigenvalue
($\lambda_3 \approx 2.9021$) that we discarded in order to obtain
a two-dimensional representation. By omitting this third dimension,
we are approximating the full geometry in a lower-dimensional space.
As a result, some distances are distorted.

\medskip

In other words, classical MDS with $p=2$ does not reproduce the
original distance matrix exactly because the data are not perfectly
two-dimensional. The two largest eigenvalues dominate, which is why
the reconstruction is reasonably accurate, but information contained
in the third eigenvalue is lost.

\medskip

The reconstructed distance matrix is therefore only approximately
equal to the original distance matrix. If we retained all positive
eigenvalues (that is, used a three-dimensional configuration),
the original distances would be recovered exactly up to numerical
precision.


